\subsection{ABSOLUTELY CONVERGENT INTEGRALS}

DEFINITION 2. One says that the integral of a regulated function \(\mathbf{f}\) over an interval \(\mathrm{I} \subset \mathbf{R}\) is absolutely convergent if the integral of the positive function \(\| \mathbf{f}(x) \|\) is convergent.

PROPOSITION 4. If the integral of \(\mathbf{f}\) over \(\mathrm{I}\) is absolutely convergent then it is convergent, and one has

\[
\left\| \int_ {I} \mathbf {f} (t) d t \right\| \leqslant \int_ {I} \| \mathbf {f} (t) \| d t.\tag{2}
\]

Indeed, for every compact interval \(J \subset I\) one has (II, p.~61, formula (16))

\[
\left\| \int_ {J} \mathbf {f} (t) d t \right\| \leqslant \int_ {J} \| \mathbf {f} (t) \| d t.\tag{3}
\]

If the integral of the positive function \(\|\mathbf{f}(x)\|\) is convergent, then for every \(\varepsilon > 0\) there exists a compact interval \([\alpha, \beta]\) contained in I, such that, for every compact interval \([x, y]\) contained in I and having no interior point in common with \([\alpha, \beta]\) , one has \(\int_{x}^{y} \|\mathbf{f}(t) dt\| \leqslant \varepsilon\) (II, p.~65, prop. 1); one deduces that \(\left\|\int_{x}^{y} \mathbf{f}(t) dt\right\| \leqslant \varepsilon\) , which shows convergence of the integral over I (II, p.~16, prop. 1); on passing to the limit in (3) one deduces the inequality (2).

COROLLARY. Let E, F, G be three complete normed spaces over R, and \((\mathbf{x}, \mathbf{y}) \mapsto [\mathbf{x}, \mathbf{y}]\) a continuous bilinear map from \(E \times F\) into G. Let f, g be two regulated functions on I, with values in E and F respectively. If f is bounded on I and if the integral of \(\mathbf{g}\) is absolutely convergent over I, then the integral of \([\mathbf{f},\mathbf{g}]\) is absolutely convergent.

Indeed, there exists a number \(h > 0\) such that one has \(\| [\mathbf{x},\mathbf{y}] \| \leqslant h \| \mathbf{x} \|\) . \(\| \mathbf{y} \|\) identically (Gen.~Top., IX, p.~173, th. 1); if one puts \(k = \sup_{x \in I} \| \mathbf{f}(x) \|\) , then one has \(\| [\mathbf{f}(x),\mathbf{g}(x)] \| \leqslant hk \| \mathbf{g}(x) \|\) on I; the comparison principle now shows that the integral of {[}f,g{]} is absolutely convergent, and, from (2),

\[
\left\| \int_ {I} [ \mathbf {f} (t). \mathbf {g} (t) ] d t \right\| \leqslant h k \int_ {I} \| \mathbf {g} (t) \| d t.
\]

Remark. An integral can be convergent without being absolutely convergent; this is what is shown by Example 3 of II, p.~64, where the series with general term \(u_{n}\) is convergent without being absolutely convergent.

